% =====================================================================
%  GMN3 §5  与簇代数的关系（Relation to cluster algebras）
%  原文：Lamination-Paper.tex 3115--3694 行；PDF p.48--56（印刷 p.47--55）
%  编号公式 (5.1)--(5.29) 共 29 条，全部在本块内，故 \eqref 可用；
%  跨块引用（节号、式号）暂写原文编号并加 % TODO，合稿时统一改 \ref/\eqref
%  注：BRIEF 分块表记本块 35 条公式，系原稿中被 % 注释掉的 4 个环境
%      （L3369 equation、L3592 align 3 行、L3642 equation、L3682 equation）
%      虚增 6 个编号单位所致；实际编译只出 29 条，已按 PDF 核对。
% =====================================================================

\section{与簇代数的关系}
\label{sec:Cluster-Algebras}

形式线算子的代数与数学中的\emph{簇代数}（cluster algebra）理论密切相关。（簇代数的文献如今已浩如烟海，下面只给出几个指路标：\cite{MR1887642,MR2132323,cluster-intro} 与 \cite{cluster-portal}。）本节简要勾勒这一关系的一小部分。我们在第 13 节% TODO \ref{sec:Open-Problems}
还会讨论，其中尚有更多值得弄清楚的事情。除了建立与数学中一个有趣课题的联系之外，我们还将看到，这一联系使我们能够只用模空间中的「局部」规则来构造形式线算子，亦即只用关于 BPS 谱的部分知识即可给出的、$\Xi$ 各腔之间的局部变换规则。此外，我们还将说明：仅凭对 BPS 谱的部分知识，原则上就足以构造出整个 BPS 谱。

\subsection{簇代数的基本定义}
\label{subsec:Cluster-Defs}

簇代数的定义从一个\emph{种子}（seed）开始：种子就是一个 $m\times n$ 矩阵 $B_{ij}$，其中 $m\geq n$、元素为整数，并且其上方 $n\times n$ 块可斜对称化。也就是说，存在正整数 $d_i$，使得 $1\leq i,j\leq n$ 时的 $B_{ij}(d_j)^{-1}$ 反对称：$d_i B_{ij} = - d_j B_{ji}$。对每个指标 $1\leq k \leq n$，沿方向 $k$ 的\emph{种子变异}（seed mutation）定义为变换 $\mu_k: B\to B'$：
%
\begin{equation}\label{eq:seed-mtrx}
B'_{ij} = \begin{cases} - B_{ij} & {\rm If} ~~ i = k ~~ {\rm or}~~ j=k \\
B_{ij} + {\rm sgn}(B_{ik}) [B_{ik} B_{kj}]_+ & {\rm Else} \\
\end{cases}
\end{equation}
%
其中我们引入了记号 $[x]_+ = {\rm Max}\{ 0,x\}$，以及
%
\begin{equation}
{\rm sgn}(x) = \begin{cases} 1 & x>0 \\  0 & x =0 \\ -1 & x < 0 \\ \end{cases}
\end{equation}
%
簇变量是 $A_1, \dots, A_m$。对每个 $k$（$1\leq k \leq n$），簇变换定义为：当 $i\not= k$ 时 $\mu_k(A_i) = A_i$，而
%
\begin{equation}
\mu_k(A_k) = \frac{ \prod_{j\vert B_{kj}>0} A_j^{B_{kj}} + \prod_{j\vert B_{kj}<0} A_j^{-B_{kj}} }{A_k}
\end{equation}
%
这一递归过程定义了某个代数的生成元序列，这个代数就称为簇代数。指标满足 $n<i\leq m$ 的变量 $A_i$ 称为「冻结变量」（frozen variables）。簇代数具有许多漂亮而出色的性质。

Fock 与 Goncharov 在 \cite{MR2567745} 中证明：若定义
%
\begin{equation}
x_i = \prod_j A_j^{B_{ij} }
\end{equation}
%
则 $x_i$ 相应的变换为 $\mu_k: x_i \to x_i'$，其中
%
\begin{equation}\label{eq:Mutation-x}
 x_i' = \begin{cases}  x_k^{-1} & i=k \\
x_i (1+ x_k^{-{\rm sgn}(B_{ik})})^{-B_{ik}} & i \not= k \\
\end{cases}
\end{equation}
%
Kontsevich 与 Soibelman \cite{ks1} 强调：式 \eqref{eq:Mutation-x} 可以用我们称之为 KS 变换的辛变换来解释。

\subsection{$\CN=2$ 理论中的簇代数结构}
\label{subsec:Seeds-from-N=2}

为了与 $\CN=2$ 理论建立联系，我们按如下方式进行。

相当一般地，若 $\Omega(\gamma;u)\not=0$，我们就称荷 $\gamma \in \Gamma_u$ 为\emph{根}。记 $\CR(u)$ 为固定 $u$ 处的根系。我们把\emph{正根系}定义为一个不交分解 $\CR(u) = \CR^+(u) \amalg \CR^-(u)$，其中若 $\gamma \in \CR^+(u)$，则 $-\gamma \in \CR^-(u)$。我们只考虑这样的正根系：其对应的 BPS 射线位于 $\zeta$ 平面内的某个半空间中。反过来，在 $\zeta$ 平面内选定一个半空间，也就决定了根系到正根与负根的一个分解。给定一个正根系，我们把\emph{单根}定义为不能写成另外两个正根之和的正根。给定任意一点 $(u,\zeta)$，我们可以典范地定义一个正根系：把满足 ${\rm Im}Z_{\gamma}/\zeta >0$ 的那些根宣布为正根。等价地，给定一个不在 BPS 墙上的点 $(u,\zeta)$，可以这样来定义正根系：取所有 BPS 射线落在过 $\zeta$ 的直线之逆时针一侧半平面内的那些根。固定 $u$ 时，点 $(u,\zeta)$ 落在 $\zeta$ 平面内的一个腔中，该腔由两道墙围成：相对于上述典范正根系，一道是逆时针一侧某个单根的墙，另一道是顺时针一侧某个「负单根」——即单根取负——的墙。若固定 $\zeta$ 而改变 $u$，则除非某条 $\pm$ 单根的 BPS 射线扫过 $\zeta$，$(u,\zeta)$ 不会穿越到 $\Xi$ 的另一个腔中去。于是，我们可以用单根系 $\Delta$ 来标记 $\Xi$ 中的各个腔。

现在，对 $\Xi$ 的任一腔 $c$，用单根 $\{\gamma_i\}$ 来定义矩阵
%
\begin{equation}
B_{ij} = \langle \gamma_i, \gamma_j\rangle   \Omega(\gamma_j;u)
\end{equation}
%
只要诸 $\Omega(\gamma_j;u)$ 全为正，这个矩阵显然可斜对称化。冻结指标 $n<i\leq m$ 标记味子格的生成元。下面我们为单根集合提出一条跨 BPS 墙的变换规则，它保证：当 $(u,\zeta)$ 沿某条路径从一个腔走到下一个腔时，发生的正是种子变异。正如刚才所观察到的，该路径必定穿过某条 $\pm$ 单根的墙。路径在 $\zeta$ 平面内的投影逆时针地越过某个单根 $Z_{\gamma_k}$，也就是从 ${\rm Im}Z_{-\gamma_k}/\zeta <0 $ 的区域进入 ${\rm Im}Z_{-\gamma_k}/\zeta > 0 $ 的区域。这就是在穿过墙 $\widehat{W(-\gamma_k)}$。类似地，若路径顺时针越过某个单根的负值 $Z_{-\gamma_k}$，从 ${\rm Im}Z_{\gamma_k}/\zeta >0 $ 的区域进入 ${\rm Im}Z_{\gamma_k}/\zeta < 0 $ 的区域，则它穿过的是墙 $\widehat{W(\gamma_k)}$。

对顺时针穿过 $\widehat{W(\gamma_k)}$ 的路径，变换
%
\begin{equation}\label{eq:charge-mutation}
\mu_{k,+}: \gamma_i \to \gamma_i' = \begin{cases} - \gamma_k & i=k \\   \gamma_i + \gamma_k [B_{ik}]_+  & i \not=k\\
\end{cases}
\end{equation}
%
给出新腔 $\mu_{k,+}(c)$ 中的新单根系。类似地，对逆时针穿过 $\widehat{W(-\gamma_k)}$ 的路径（其中 $\gamma_k$ 是单根），变换
%
\begin{equation}
\mu_{k,-}: \gamma_i \to \gamma_i' = \begin{cases} - \gamma_k & i=k \\   \gamma_i - \gamma_k [B_{ik}]_-  & i \not=k\\
\end{cases}
\end{equation}
%
给出新腔 $\mu_{k,-}(c)$ 中的新单根系。这里我们引入了记号 $[x]_- := {\rm Min}\{x,0\}$。可以验证，$\mu_{k,+}$ 是 $\mu_{k,-}$ 的双侧逆。

现在可以验证：只要 $\Omega(\gamma_j';u) = \Omega(\gamma_j;u)$，新腔的矩阵 $B'_{ij}$ 与旧腔的矩阵 $B_{ij}$ 之间的关系（对 $\mu_{k,\pm}$ 两者都是如此）恰好就是式 \eqref{eq:seed-mtrx} 所定义的种子变异 $\mu_k$。

此外，Kontsevich--Soibelman 变换
%
\begin{equation}
\CK_{\gamma_0}^{-\Omega(\gamma_0;u)} x_\gamma = x_{\gamma} (1-\sigma(\gamma_0) x_{\gamma_0})^{-\langle \gamma,\gamma_0\rangle \Omega(\gamma_0)}
\end{equation}
%
与簇变换 \eqref{eq:Mutation-x} 相合，前提是取 $\sigma(\gamma_0)=-1$、$x_i' = x'_{\gamma_i'}$、$x_i = x_{\gamma_i}$、$\gamma_0 = \gamma_k$ 以及 $\gamma_i' = \mu_{k,+}\gamma_i$。因此，举例来说，若定义 $\hat x_{\gamma} = \sigma(\gamma) x_{\gamma}$，则带框 BPS 简并数可以按照式 (3.35)% TODO \eqref{eq:Careful-rule}
用单根变异所诱导的簇变换来计算。

借助量子簇代数理论 \cite{MR2233852}，这一陈述可以推广到完整的 $y \neq \pm 1$ 穿墙；但为简单起见，我们只看经典情形。

若跨腔的变换满足上述性质，我们就说一个 $\CN=2$ 场论具有\emph{簇代数性质}。也就是说：对单根有 $\Omega(\gamma_i;u)>0$，并且 $\Omega(\mu_{k,\pm} \gamma_j;u) = \Omega(\gamma_j;u)$。（为避免对这个第二条件产生误解，请注意：腔壁\emph{不是}边缘稳定墙，因此该条件是良定的。）

具有簇代数性质的主要好处之一是：可以给出一个构造理论中 BPS 谱与线算子的算法。为看清这一点，先考虑 BPS 谱。假定我们知道某个腔中各单根的 BPS 简并数，并假定知道 $Z_{\gamma_i}(u)$ 的取值。这些量在腔内处处变化，但始终留在某个半平面内。固定 $u$，在该腔内取一个 $\zeta$ 值（于是 ${\rm Im}Z_{\gamma_i}(u)/\zeta >0$）。现在让 $\zeta$ 沿逆时针方向移动，它会撞上某个单根 $\gamma_k$ 的 BPS 谱线。此时便可借助变异 $\mu_{k,-}$ 重新算出新的单根基。继续逆时针转动 $\zeta$，就可以这样一路走下去。若能一直把 $\zeta$ 转过整整 $\pi$ 的角度，则整个 BPS 谱就构造出来了。正如我们将在下面两个例子中看到的，这样做可能遇到阻碍。基本的原因是：所有腔的单根系都只含超多重态，而更高自旋粒子的 BPS 射线被无穷多条、属于无穷多个超多重态的射线所包围，这些射线向高自旋射线聚积。因此高自旋射线从不与某个腔的边界相联系。若只在 $\zeta$ 平面内平移，过程会在这一无穷超多重态序列的第一个上就卡住，无法到达序列「末端」的高自旋粒子。另一方面，在 $A_1$ 理论中，并且可能在 $\CS$ 类中的所有理论中，看来只要允许在 $(Z_\gamma, \zeta)$ 的空间中沿一般路径移动，所有腔都能经由有限多步连通。于是，一个可能的谱计算策略是：利用这类路径「跳」到阻碍的前方，然后再回溯，以控制阻碍两侧的无穷超多重态序列，从而求出跨过阻碍的剩余坐标变换。该变换还可进一步分解为变换 $\CK_\gamma$，由此读出理论的完整谱。

作为上述步骤的一个例子，设 $\Gamma$ 是二维的、$\Omega(\gamma_1)=\Omega(\gamma_2)=1$，并且对某个正整数 $n$ 有
%
\begin{equation}
B_{ij} = \begin{pmatrix} 0 & n \\ -n & 0 \\ \end{pmatrix}.
\end{equation}
%
从 $\zeta$ 落在实轴上、且 ${\rm Im} Z_1 > 0 $ 与 ${\rm Im} Z_2 >0$ 的那个腔出发。若 ${\rm arg}Z_1 > {\rm arg} Z_2$，我们就作一列变换 $\mu_{2,-}, \mu_{1,-}, \mu_{2,-}, \cdots$。所得单根的图样周期为 4，具体就是
%
\begin{equation}
\begin{split}
\gamma_1^{[1]}=   \gamma_1, & \quad   \gamma_2^{[1]}= - \gamma_2, \\
\gamma_1^{[2]}= - \gamma_1, & \quad \gamma_2^{[2]}= -\gamma_2, \\
\gamma_1^{[3]}=  - \gamma_1, & \quad \gamma_2^{[3]}=  \gamma_2.
\end{split}
\end{equation}
%
反之，若 $\arg Z_2 > \arg Z_1$，则逆时针移动 $\zeta$ 时会遇到变异序列 $\mu_{1,-}, \mu_{2,-}, \mu_{1,-},\cdots$。若把初始的腔标为 $j=0$，我们得到一列腔，满足
%
\begin{equation}
\begin{split}
\gamma_1^{[2j+1]} & = -\gamma_1^{[2j]}  , \\
\gamma_2^{[2j+1]} & = n\gamma_1^{[2j]}+ \gamma_2^{[2j]},
\end{split}
\end{equation}
%
\begin{equation}
\begin{split}
\gamma_1^{[2j+2]} & =\gamma_1^{[2j+1]}+ n\gamma_2^{[2j+1]},   \\
\gamma_2^{[2j+2]} & =  - \gamma_2^{[2j+1]}.
\end{split}
\end{equation}
%
当 $n=1$ 时，这一列变换的周期为 $6$。它已在 $N=3$ Argyres-Douglas 理论的弱耦合区中被具体实现，见第 4.3 节% TODO \ref{sec:formaln3}
。当 $n=2$ 时，这列变换没有有限的周期性，且 $\mu_{2,-}\mu_{1,-}$ 是由矩阵
%
\begin{equation}
1 + N = \begin{pmatrix} 3 & 2 \\ -2 & -1 \\ \end{pmatrix}
\end{equation}
%
所给出的变换。由于 $N^2 =0$，可以立刻写出
%
\begin{equation}
\begin{split}
\gamma^{[2j]}_1 & = (2j+1)\gamma_1 + 2j  \gamma_2, \\
\gamma^{[2j]}_2 & =  -2j \gamma_1  -  (2j-1) \gamma_2.
\end{split}
\end{equation}
%
这在 $SU(2)$ $N_f=0$ 理论中实现过，见第 4.4 节% TODO \ref{subsec:formal-SU2}
。注意这些 BPS 射线聚积在沿 $\pm (Z_1 + Z_2)$ 的射线上。因此，无法越过这个极限继续转过整整 $\pi$ 的角度。不过，若改沿顺时针方向继续，就会找到对应于相反射线的第二个极限值，于是实际上完整的谱仍然被生成出来。

当 $n>2$ 时会出现一个有趣的现象。从 $2j$ 到 $2j+2$ 的变换矩阵是
%
\begin{equation}
\begin{pmatrix} n^2 -1 & n \\ -n &  -1 \\ \end{pmatrix} = S^{-1} \begin{pmatrix} \lambda_+ & 0 \\
0 & \lambda_- \\ \end{pmatrix} S
\end{equation}
%
其中 $1+\lambda_\pm = \half (n^2  \pm n \sqrt{n^2-4} ) $。沿逆时针方向移动时，BPS 射线的极限斜率为
%
\begin{equation}\label{eq:slope-ccw}
\frac{ \xi {\rm Im}Z_1 +   {\rm Im}Z_2 }{\xi {\rm Re} Z_1 +   {\rm Re} Z_2}
\end{equation}
%
其中 $\xi = \frac{\lambda_+ +1 }{n} > 1$。现在若考虑沿顺时针方向转动，我们会遇到变异序列 $\mu_{2,+}$, $\mu_{1,+}$, $\mu_{2,+},\dots$，它给出一列单根
%
\begin{equation}
\begin{split}
\gamma_1^{[2j+1]} & =   \gamma_1^{[2j]} + n   \gamma_2^{[2j]} \\
\gamma_2^{[2j+1]} & =  -   \gamma_2^{[2j]} \\
\end{split}
\end{equation}
%
\begin{equation}
\begin{split}
\gamma_1^{[2j+2]} & =     -\gamma_1^{[2j+1]}    \\
\gamma_2^{[2j+2]} & = n \gamma_1^{[2j+1]} + \gamma_2^{[2j+1]}  \\
\end{split}
\end{equation}
%
此时的极限斜率为
%
\begin{equation}\label{eq:slope-cw}
\frac{ {\rm Im}Z_1 +  \xi  {\rm Im}Z_2 }{  {\rm Re} Z_1 +  \xi {\rm Re} Z_2}
\end{equation}
%
而该射线指向 $-\pi < {\rm arg}\zeta < -\pi/2$ 的那个象限。由于 $\xi >1$，两条极限 BPS 射线之间张开的角小于 $\pi$。因此存在一个「空隙区域」，如图所示~\ref{fig:gap}，上述步骤无法探查到那里。该区域内被占据的荷的谱已经有人算出 \cite{gctv}，其结果清楚地显示了困难所在：在该区域中，$\Omega(\gamma) \neq 0$ 的那些 $Z_\gamma$ 的相位是\emph{稠密}的，因此谈论腔其实已不太可能。

\insfig[0.58\textwidth]{gap}{在秩 2 例子中当 $n>2$ 时，我们只能用固定 $u$ 处的簇变换来确定落在某个角小于 $\pi$ 的角扇区内的那些射线的 BPS 谱，该扇区用蓝色箭头标出。图中对 $Z_1$ 与 $Z_2$ 的相位作了一个方便的选取。}

当一个 $\CN=2$ 理论具有簇代数性质时，会立刻得到一些好处。文献 \cite{MR1888840} 的洛朗现象使我们能够构造形式叶层。形式线算子的构造涉及某种「量子洛朗现象」，它似乎尚未在数学文献中被讨论过（尽管它与 Fock 与 Goncharov 的「泛洛朗多项式」相去不远）。若我们的目的只是检验某个形式线算子的强正性，那么其实并不关心荷格的标号；于是只需把生成函数用 $X_i$ 表出，忘掉 $\gamma_i$，并确保强正性被所有可能的簇变换保持即可。这可以用来简化我们在第 4.3 节% TODO \ref{sec:formaln3}
与第 4.4 节% TODO \ref{subsec:formal-SU2}
中给出的推导，并轻而易举地把它们推广到弱耦合区，因为强耦合与弱耦合处的簇变换序列无从区分。例如，若忘掉荷的标号，$N=3$ Argyres-Douglas 理论总共只有五个不同的腔，由五边形的五种三角剖分标记。簇变换的周期为 $5$，而基本的形式线算子只是经历一个周期序列 $X\to X+X \to X+X+X \to X+X \to X \to X$。我们在强耦合处遇到的周期为 $20=4 \times 5$ 的长序列、以及在弱耦合处周期为 $30 = 6 \times 5$ 的长序列，不过是周期 $5$ 的簇变换与周期 $4$ 或 $6$ 的荷重标号的组合而已。

我们猜想：$\CS$ 类中的理论都具有簇代数性质。这一猜想受到 Fock 与 Goncharov 关于高阶 Teichm\"uller 理论的工作 \cite{MR2233852} 的启发。若我们的猜想成立，则上述算法或许提供了计算高秩 $\CN=2$ 规范理论 BPS 谱的一条途径。我们认为这是一个重要的想法，值得继续追下去。基于以上观察，一个自然会问的问题是：是否每个簇代数都源自某个 $\CN=2$ 理论？这个问题我们也留待将来。

\subsection{$A_1$ 理论中的簇代数结构}
% 原文此小节的 label 与 5.2 重复（都写作 subsec:Seeds-from-N=2），系原文笔误；此处加后缀区分。
\label{subsec:Seeds-from-A1}

本节我们将证明：$A_1$ 型理论这一子类确实具有簇代数性质。这一点可以借助 \cite{Gaiotto:2009hg} 中所解释的腔与 $C$ 的 WKB 三角剖分之间的关系来证明。回忆该文第 6 节：给定 $(u,\zeta)$，$C$ 上有一个典范的（WKB）三角剖分。该文第 7.8 节证明了还存在一个与之相伴的单根系 $\{ \gamma_E \}_{E\in \CE(T)}$。而且，每个这样的单根上都承载着单个超多重态，其 PSC 等于 $1$。当 $(u,\zeta)$ 变动时，除非某条边发生 flip，三角剖分只作同痕变化；而后者只可能在 $(u,\zeta)$ 穿过一道 BPS 墙时发生。因此，腔的墙是诸墙 $\widehat{W}(\pm \gamma_E)$（$E\in \CE(T)$）的一个子集。一般地，我们预期所有边都会贡献腔的墙。由一条边 $E_k$ 的 flip 所引起的 $\{ \gamma_E \}_{E\in \CE(T)}$ 的变换规则，已在 \cite{Gaiotto:2009hg} 的第 7.6 节中算出，它们与 $\mu_{k,\pm}$ 相符。

我们可以用 $C$ 的理想三角剖分，连同某个「适配二重覆盖」的拓扑类，来标记各个腔。$C$ 的一个适配二重覆盖，是指在 $C$ 的每个三角形中各有一个分支点的二重覆盖，并且要选定各边 $E$ 之提升的定向，使得在 $C$ 的奇点的原像处，这些提升要么全是入的、要么全是出的。

\subsection{$\CN=2^*$ $SU(2)$ 规范理论中的形式线算子}
\label{subsec:Formal-N2-star}

现在我们转向规范代数为 $su(2)$ 的 $\CN=2^*$ 理论。它由 $\CN=4$ 理论给一个伴随超多重态赋予质量 $m$ 而得到。它就是 $C$ 为带一个穿刺点的环面时的 $A_1$ 理论。\footnote{更确切地说，$C$ 为带穿刺点环面时的 $A_1$ 理论，是 $\N=2^*$ 理论再加上一些解耦的场，即在该穿刺点对应的 $SU(2)$ 味对称下的一个半超多重态二重态。}Seiberg--Witten 微分由
%
\begin{equation}
\lambda^2 = \left( m^2 \wp(z\vert \tau) + u\right) (dz)^2,
\end{equation}
%
给出，其中 $z$ 是模数为 $\tau$ 的椭圆曲线上的平坦坐标，$\wp(z\vert \tau)$ 是 Weierstrass 函数。Seiberg--Witten 曲线亏格为 2、带两个穿刺点。局部系统 $\Gamma$ 的秩为 $3$，可以取生成元 $\gamma_1, \gamma_2, \gamma_3$，使得对 $i=1,2,3$ 有 $\langle \gamma_i, \gamma_{i+1} \rangle =2$，其中指标 $i$ 按 $3$ 阶循环理解。注意 $\gamma_1 + \gamma_2 + \gamma_3$ 生成 $\langle \cdot, \cdot \rangle$ 的一维零化子。这就是秩 $1$ 的「味格」。

可归一化参数是 $u$，并且 $\CB^* = \IC-\CS$，其中 $\CS$ 由三个奇点 $u_i$ 组成，满足 $\wp(z_i\vert \tau) = - u_i/m^2$ 与 $\wp'(z_i\vert\tau) =0$。局部系统 $\Gamma$ 绕这三个点有单值变换。众所周知该理论具有 $SL(2,\IZ)$ 对称性，其中 $PSL(2,\IZ)$ 以标准方式作用在耦合常数 $\tau$ 上。

遗憾的是，$\CN=2^*$ 理论的 BPS 谱 $\Omega(\gamma;u)$ 并不知道其显式形式。\footnote{我们注意到，该谱隐含地编码在 \cite{Gaiotto:2009hg} 的「谱生成元」之中。这个变换容易写出，但要把它分解为 Kontsevich--Soibelman 因子却并不直接。}尽管如此，很清楚该谱与迄今考察过的例子有一些鲜明的差别。特别是，这里没有谱有限的「强耦合区」。我们预期谱总是无穷的，并由数目可变的若干矢量多重态控制，而每个矢量多重态都伴随着无穷多个超多重态。因此腔的结构非常复杂。于是，这里正是一个可以非常有效地使用局部簇变换规则的例子。

在每个腔 $[c]$ 中有三个单根 $\gamma_i[c]$；该腔由六道墙 $\widehat{W}(\pm \gamma_i[c])$ 围成。$\sum_i \gamma_i[c]$ 是纯味的，并且我们\emph{或者}有 $\langle \gamma_i[c],\gamma_{i+1}[c]\rangle=2$，\emph{或者}有 $\langle \gamma_i[c],\gamma_{i+1}[c]\rangle=-2$。我们分别称这些腔为正型与负型。该腔就是集合 $\{ (u,\zeta): {\rm Im}Z_{\gamma_i[c]}(u)/\zeta > 0 \}$。注意一次变异 $\mu_{i,\pm}$ 把正型的腔变到负型的腔。

方便起见，定义
%
\begin{equation}
[p,q,r]:= X_{\half (p \gamma_1 + q \gamma_2 + r \gamma_3)}.
\end{equation}
%
当需要强调腔时，我们可以写成 $[p,q,r](c)$。我们首先注意到，与先前的例子相反，这里没有任何单项式 $[p,q,r]$ 能够生成形式线算子。在我们的腔中，${\rm Im} Z_i/\zeta >0$（$i=1,2,3$）。检视规则 (4.3)% TODO \eqref{eq:explicit-dilog-conj}
与 (4.4)% TODO \eqref{eq:explicit-dilog-conj-ii}
，我们看到必须有 $\langle \gamma, \gamma_i[c] \rangle <0$（$i=1,2,3$），其中 $\gamma = \half (p\gamma_1[c]+q\gamma_2[c]+r\gamma_3[c])$。另一方面，$\gamma_1[c]+\gamma_2[c]+\gamma_3[c]$ 位于 $\langle \cdot, \cdot \rangle$ 的核中，因此这三个不等式永远不可能同时满足。

尽管无法由单项式生成形式线算子，却有三个漂亮的类比物，它们对应于 Wilson 线算子 $W$（后者在第 4.4 节% TODO \ref{subsec:formal-SU2}
中十分有用）：
%
\begin{equation}\label{eq:star-formal-Wilson}
\begin{split}
W_1[c] & = [0,1,1] + [0, -1, -1] + [0,-1,1],\\
W_2[c] & = [1,0,1] + [-1,0, -1] + [1,0,-1],\\
W_3[c] & = [1,1,0] + [-1, -1, 0] + [-1,1,0].
\end{split}
\end{equation}
%
若腔是正型的（$\langle \gamma_i[c],\gamma_{i+1}[c]\rangle=+2$），这些就是跨过该腔所有墙都变换良好的唯一组合。对负型的腔，$\langle \gamma_i[c],\gamma_{i+1}[c]\rangle=-2$，这三个 Wilson 算子具有如下形式
%
\begin{equation}
\begin{split}
W_1[c] & = [0,1,1] + [0, -1, -1] + [0,1,-1],\\
W_2[c] & = [1,0,1] + [-1,0, -1] + [-1,0,1],\\
W_3[c] & = [1,1,0] + [-1, -1, 0] + [1,-1,0].
\end{split}
\end{equation}
%
记 $\mu_{i,\pm}$ 为跨过墙 $\widehat{W}(\pm \gamma_i[c])$ 的变换（变异）。当 $j\not=i$ 时，$W_i[c]$ 在 $\mu_{j,\pm}$ 下的变换遵循 $SU(2)$ Wilson 线那个熟悉的模式：一个二项（自旋 $1/2$）的 halo 塌缩，而某一项长大成一个二项的 halo，于是我们在下一个腔中仍然得到一个三项的表达式。也就是说，对 $i \not=j$ 我们有简单的结果
%
\begin{equation}
\mu_{i,\pm}(W_j[c]) = W_j[\mu_{i,\pm}(c)].
\end{equation}
%
另一方面，若用 $\mu_{i,\pm}$ 去变异 $W_i[c]$，则有两项不变，而第三项变成一个三项的（自旋 $1$）halo。例如，用 $\mu_{3,+}$ 变换到腔 $c' = \mu_{3,+}(c)$，我们得到：
%
\begin{equation}\label{eq:mu-three-W3}
\begin{split}
\mu_{3,+} W_3[c] &=  [1,1,0](c) + [-1, -1, 0](c) + ([-1,1,0]+ [2] [-1,1,2] + [-1,1,4])(c)\\
& = [1,1,2](c') + [-1,-1,-2](c') + \left( [-1,1,2] + [2] [-1,1,0] + [-1,1,-2]\right)(c'). \\
\end{split}
\end{equation}
%
与 $SU(2)$ $N_f=0$ 的情形一样，Wilson 线算子的环关系与它们跨腔壁的变换性质美妙地联系在一起。例如
\begin{align}
W_1[c] W_2[c]&= y^{-\half}W_3[c]+y^{\half}\bigl(X_{\half \gamma_1[c]+\half \gamma_2[c]+\gamma_3[c]} +X_{\half \gamma_1[c]-\half \gamma_2[c]+\gamma_3[c]} \notag\\
&\quad +X_{-\half \gamma_1[c]-\half \gamma_2[c]-\gamma_3[c]} +X_{\half \gamma_1[c]-\half \gamma_2[c]-\gamma_3[c]} + (y+y^{-1})X_{-\half \gamma_2[c]+\half\gamma_1[c]}\bigr)
\end{align}
乘以 $y^\half$ 之后的生成函数在跨过任一 $\pm \gamma_3[c]$ 墙时会简化：它失去一个 Fock 空间，变成一个简单的 Wilson 线，这一点从 \eqref{eq:mu-three-W3} 可以认出。

这类乘积的一般关系是：
%
\begin{equation}\label{eq:quant-skein-rel}
W_{i}[c]  W_{i+1}[c] = y^{-1/2} W_{i+2}[c]   + y^{1/2} \mu_{i+2,+} W_{i+2}[\mu_{i+2,-}(c)]
\end{equation}
%
对 $W_{i+1}[c] W_{i}[c]$，只需把 $y\to 1/y$。

我们预期，通过逐次乘以 Wilson 线，可以递推地生成整个形式线算子环，正如第 4.4 节% TODO \ref{subsec:formal-SU2}
的情形那样。

我们正一步步揭示该理论线算子中一个美妙的结构。在本文的后半部分，我们将发展更有力的工具来分析所有 $A_1$ 理论；在此过程中，我们会重新发现本节的这些公式。为了与第 8.4 节% TODO \ref{subsec:N2-star-moduli}
和第 10.7 节% TODO \ref{subsec:SU2-N2star-Examples}
作比较，我们再作最后一个计算：注意 $[2,2,2]$ 与 $[-2,-2,-2]$ 是中心元，并且可以适当地选取系数 $a_{ijk}\in \IZ[y^{1/2},y^{-1/2}]$，使得
%
\begin{equation}\label{eq:quant-mod-eq}
\begin{split}
W_1^2 + W_2^2 + W_3^2 - a_{ijk} W_i W_j W_k &= 1 + y^{-2} + B(y^2-y^{-2}) \\
&\quad - \left(y^{-1} + (y-y^{-1})B\right)([2,2,2]+[-2,-2,-2])
\end{split}
\end{equation}
%
是中心算子。这里 $B$ 是任意常数。式 \eqref{eq:quant-mod-eq} 的 $y \to 1$ 极限就是式 (8.9)% TODO \eqref{eq:moduli-equation}
。
